\documentclass[10pt,a4paper]{amsart}
\usepackage[margin=19mm]{geometry}
\usepackage{amsmath,amssymb,mathtools}
\usepackage[scr=rsfs,cal=euler]{mathalfa}
\usepackage{xfrac}
\usepackage[unicode,hidelinks,bookmarks=false]{hyperref}
\newcommand{\ie}{i.e.\ }
\newcommand{\cf}{cf.\ }
\newcommand{\resp}{respectively\ }
\newcommand{\Subsection}{\subsection}
\providecommand{\nobreakdash}{\nobreak\hbox{-}}
\setlength{\emergencystretch}{2em}
\multlinegap0pt
\usepackage[shortlabels]{enumitem}
\usepackage{xcolor}
\newtheorem{definition}{Definition}
\newtheorem{situation}{Situation}
\newtheorem{theorem}{Theorem}
\renewcommand{\geq}{\geqslant}
\renewcommand{\leq}{\leqslant}
\renewcommand{\mathbb}{\mathbf}
\renewcommand{\setminus}{\mathbin{\rule[0.2em]{0.67em}{0.12em}}}%
\renewcommand{\subset}{\subseteq}
\title{Asymptotic Betti bounds for hypersurfaces in a singular variety}
\author{Xuanyu Pan, Dingxin Zhang, Xiping Zhang}
\date{Source: arXiv:2601.20179v1 (CC BY 4.0)}
\begin{document}
\maketitle

\noindent\emph{Source and licence.} Asymptotic Betti bounds for hypersurfaces in a singular variety. Version arXiv:2601.20179v1 (CC BY 4.0), distributed by arXiv under the Creative Commons Attribution 4.0 International licence, http://creativecommons.org/licenses/by/4.0/, as recorded in the arXiv licence metadata for this exact version and shown on the abstract page https://arxiv.org/abs/2601.20179v1. Full licence notice: \texttt{licenses/arxiv-2601.20179-CC-BY-4.0.txt}.\par\smallskip

\noindent\textit{Editorial scope.} Counted unit: the article's theorem theorem:arbitrary (source0.tex lines 2141--2152). The article's complete original proof of it is delivered with it (source0.tex lines 2154--2208, one \texttt{proof} environment of 55 lines). No mathematics is written, repaired or re-derived: the statement and the proof are the article's own lines, in the article's own notation, and no retained line is substituted or re-flowed.  Retained uncounted as the source-local context the proof needs: lines 1174--1184; lines 1603--1614; lines 1640--1663.  Nothing was omitted from any retained span, so the package carries no omission record.  No retained line contains a citation, so the wrapper carries no bibliography and no bibliography entry.  No retained line contains a figure environment, an image-inclusion command or drawing code, so no image is delivered and the package declares no asset dependency.  Editorial formatting only, in addition: an AMS \texttt{amsart} preamble that restates the article's own declarations and macros used by the retained lines, namely the environments \texttt{definition}, \texttt{situation}, \texttt{theorem} on separate counters, together with the standard packages those lines need.  The retained environments print in this document as Definition 1, Situation 1, Theorem 1 in order of appearance, whereas the article prints the same items with the numbering of its own full document.  The complete native source of the article as distributed in the arXiv e-print for this exact version is bundled unchanged as \texttt{sources/193507/source0.tex}.
\begin{definition}\label{definition:vir-gen-rank}
Let \(S\) be a variety over \(k\) of dimension \(n\), and let
\(\mathcal{G} \in D^{b}_{c}(S,\Lambda)\) be a constructible complex of
\(\Lambda\)-sheaves.  We define the \emph{virtual generic rank} of
\(\mathcal{G}\) to be
\[
\mathrm{rank}^{\circ}(\mathcal{G}) = \sum_{j\in \mathbb{Z}} (-1)^{j} \operatorname{rank}_{\Lambda} \mathcal{H}^{j}(\mathcal{G}|_{S^{\circ}})
\]
where \(S^{\circ}\) is the largest purely \(n\)-dimensional Zariski
open dense subset of \(S\) over which \(\mathcal{G}\) is lisse.
\end{definition}

\section{Hypersurface restriction of perverse sheaves }
\label{sec:main-theorem}
\begin{situation}[Setup]\label{setup}
Throughout this section, let \(\mathbb{P}\) be a smooth projective
variety over \(k\) endowed with a very ample invertible sheaf
\(\mathscr{L}\).  For simplicity, we shall refer to the hypersurface
cut out by a section of \(\mathscr{L}^{\otimes d}\) as a ``degree
\(d\) hypersurface'' in \(\mathbb{P}\).  If
\(\mathcal{F} \in D^b_c(\mathbb{P},\Lambda)\), we write
\(\mathrm{cc}(\mathcal{F}) \in \mathrm{CH}_{\ast}(\mathbb{P})\)
instead of \(\mathrm{cc}_{\mathbb{P}}(\mathcal{F})\) and \(\mathrm{cc}_j(\mathcal{F})\) for \(\mathrm{cc}_{\mathbb{P},j}(\mathcal{F})\) where \(j \in \mathbb{Z}_{\geq0}\).
\end{situation}

\begin{theorem}\label{theorem:main}
Let notation be as in \ref{setup}.  Suppose \(\mathcal{P}\) is a
perverse sheaf on \(\mathbb{P}\).  Then there exist constants
\(C_1, C_2, C_3 > 0\), independent of the choice of \(Y\) below, such
that
\begin{enumerate}
\item For any degree \(d\) hypersurface \(Y \subset \mathbb{P}\),
we have
\[
\dim \mathrm{H}^{-1}(Y;\mathcal{P}|_Y) \leq r(\mathcal{P}) \cdot \delta(\mathcal{P}) \cdot d^{n(\mathcal{P})} +
C_1 \cdot d^{n(\mathcal{P})-1}.
\]
\item For any degree \(d\) hypersurface \(Y \subset \mathbb{P}\) such
that \(\mathcal{P}|_Y[-1]\) is a perverse sheaf, we have
\[
B(Y, \mathcal{P}|_{Y}) \leq r(\mathcal{P}) \cdot \delta(\mathcal{P}) \cdot d^{n(\mathcal{P})} + C_2 \cdot d^{n(\mathcal{P})-1}.
\]
\item For any degree \(d\) hypersurface \(Y \subset \mathbb{P}\) whatsoever, we
have
\[
B(Y, \mathcal{P}|_{Y}) \leq 3\cdot r(\mathcal{P}) \cdot \delta(\mathcal{P}) \cdot d^{n(\mathcal{P})} + C_3 \cdot d^{n(\mathcal{P})-1}.
\]
\end{enumerate}
\end{theorem}

\begin{theorem}\label{theorem:arbitrary}
Suppose \(\dim X = n\), and suppose \(\mathcal{F}\) is
a plain constructible sheaf on \(X\).  Then there exist a constant \(C> 0\) such
that for any degree \(d\) hypersurface \(Y \subset X\) (\(C\) is
independent of \(Y\)), we have
\[
B(Y;\mathcal{F}|_{Y}) \leq 3 \operatorname{rank}^{\circ}(\mathcal{F}) \cdot \deg(X) \cdot d^n + C \cdot d^{n-1}
\]
(for the definition of the generic rank
\(\operatorname{rank}^{\circ}\), see
Definition~\ref{definition:vir-gen-rank}).
\end{theorem}

\begin{proof}
We proceed by an induction on the dimension of \(X\).  When the
dimension is zero, there is nothing to prove.  Suppose for any polarized variety
\((D,\mathscr{L})\) of dimension \(\leq n-1\), and any constructible sheaf \(\mathcal{F}\)
on \(D\), we have found a constant \(C_{\mathscr{L}}(D,\mathcal{F})\) which makes
the inequality holds.

Let
\[
\iota\colon X \to \mathbb{P} = \operatorname{Proj}\mathrm{Sym}^{\bullet}\mathrm{H}^0(X;\mathscr{L})^{\vee}
\]
be the projective embedding induced by \(\mathscr{L}\).  Let
\(j\colon X^{\circ} \to X\) be an open immersion of satisfying the
following properties:
\begin{itemize}
\item \(X^{\circ}\) is smooth,
\item \(X^{\circ}\) contains are the generic points of the
\(n\)-dimensional irreducible components of \(X\),
\item \(\mathcal{F}|_{X^{\circ}}\) is lisse of rank equal to
\(\mathrm{rank}^{\circ}(\mathcal{F})\).
\end{itemize}
By shrinking \(X^{\circ}\), we may further assume that
\begin{itemize}
\item \(X \setminus X^{\circ}\) is a Cartier divisor.
\end{itemize}
Let \(i \colon D \to X\) be the inclusion of
the complement \(D = X \setminus X^{\circ}\) of \(X^{\circ}\).  Then
we have a distinguished triangle
\[
j_! j^{\ast}\mathcal{F}[n] \to \mathcal{F}[n] \to i_{\ast} i^{\ast} \mathcal{F}[n].
\]
Since \(D\) is a Cartier divisor, \(j\) is quasi-finite and affine.
Therefore, \(j_!j^{\ast}\mathcal{F}[n]\), hence \(\iota_{\ast} (j_!j^{\ast}\mathcal{F}[n])\), is perverse.

Let \(Y\) be any degree \(d\) hypersurface in \(X\).
Let \(H \subset \mathbb{P}\) be the degree \(d\) hypersurface in \(\mathbb{P}\) such that \(Y = X \cap H\).
Then by
Theorem~\ref{theorem:main}(3) that we have
\[
B(Y, j_!j^{\ast}\mathcal{F}[n]|_Y) = B(H, \iota_{\ast} (j_!j^{\ast}\mathcal{F}[n])|_{H})
\leq 3\operatorname{rank}^{\circ}(\mathcal{F}) \cdot \deg(\overline{X}{}^{\circ})
\cdot d^n + C_3 \cdot d^{n-1}
\]
for some constant \(C_3\).  Here, \(\overline{X}{}^{\circ}\) is the
closure of \(X^{\circ}\).  The result then follows from an induction
on dimension: because
\[
B(Y, \mathcal{F}|_{Y}) \leq B(Y, j_!j^{\ast}\mathcal{F}[n]|_Y) + B(Y \cap D, i^{\ast} \mathcal{F}|_{Y\cap D}),
\]
we can simply take
\[
C = C_{\mathscr{L}}(X,\mathcal{F})=C_{3} + \max\left\{3\operatorname{rank}^{\circ}(i^{\ast}\mathcal{F})\cdot \deg(D), C_{i^{\ast}\mathscr{L}}(D,i^{\ast}\mathcal{F})\right\}.
\]
This completes the proof.
\end{proof}

\end{document}
